\section{A growing two-power family in the same arithmetic window}
\label{sec:two-power-uniform-actual-window}

\paragraph{Status and literal input.}
This conditional-interface proof applies the original-coefficient strict
vanishing theorem AC-G173 at coefficient $269$.
It retains every omission of AC-G171-window and adds a family whose
number of finite characters grows with the actual two-power page.
The strict theorem has no automorphy, endpoint, support or regulator
premise.  The separate common abelian window theorem AC-G172 is not an
input here, and its stronger source package is not added to the gate.
All unlisted terms and the mixed realization remain open.

Put $p=269$, $K=\mathbb Q(\omega)$, $t=\sqrt[3]p>0$, and let
$\chi$ describe $K(t)/K$.  For the retained datum, a fixed
$0<\varepsilon<1/256$, the original page and each prescribed pair $q,d$,
keep exactly
\[
 \begin{aligned}
 F_u&=\mathbb Q(\zeta_{2^{h(u)+2}},\zeta_7,
                                  \sqrt{-q},\sqrt{-d},\omega),\\
 H_u&=F_u(t),\qquad G_u=\operatorname{Gal}(H_u/K),\qquad D_u=|G_u|,\\
 u^{2+8\varepsilon}&\leq2^{h(u)}\leq2u^{2+8\varepsilon},
                      \qquad T=\lfloor u^{399}\rfloor .
 \end{aligned}
 \tag{TPW.1}
\]
For $u\geq2$ we have $h(u)\geq2$.
The coefficient rings $O_\ell$, fields $E_\ell$, ranks $b_\ell$,
strict conditions and derived indices are exactly (FC.1)--(FC.2),
(FL.2), (EU.2) and (TC.1).  Every coefficient prime is retained.

\paragraph{The complete previously omitted set.}
Write $\delta_a=\chi_a|_{G_K}$ for $a=-4,8,-8$.
Let $\phi$ be the odd primitive quartic character modulo sixteen with
$\phi(13)=i$ and kernel $\{1,7\}$, and put
$\theta_{\sigma,b}=(\phi^\sigma\chi_{-4}^{\,b})|_{G_K}$,
where $\sigma=\pm1$ and $b=0,1$.
Define
\[
 \begin{aligned}
 \mathcal E&=\{89,92,99,168,175,178\},\\
 \mathcal F_{1,0}&=\{243,248\},&\mathcal F_{-1,0}&=\{19,24\},\\
 \mathcal F_{1,1}&=\{11,60,229\},&\mathcal F_{-1,1}&=\{38,207,256\}.
 \end{aligned}
 \tag{TPW.2}
\]
The residues in these sets are modulo $268$.
The set $Z^{\rm qref}$ of AC-G171-window is the disjoint union below;
in every row $1\leq e<T$ is required:
\[
 \begin{array}{c|l|l}
 \ell&\text{finite label}&\text{weight condition}\\ \hline
 p&\chi^{-j}\phi^{-\sigma},\ j=1,2,\ \sigma=\pm1
                    &e\equiv-1\pmod {67}\\
 p&\chi^{-j}\delta_a,\ j=1,2,\ a=-4,8,-8
                    &e\equiv-1\pmod {67}\\
 p&\chi^j,\ j=1,2&e\equiv0\pmod {268}\ \text{or }e\equiv-1\pmod {67}\\
 p&1&\text{all }e\\
 p&\delta_{-4},\delta_8&e\not\equiv0\pmod {268}\\
 p&\delta_{-8}&e\bmod268\notin\{0\}\cup\mathcal E\\
 p&\theta_{\sigma,b},\ \sigma=\pm1,\ b=0,1
                    &e\bmod268\notin\mathcal F_{\sigma,b}\\
 3&1&\text{all }e.
 \end{array}
 \tag{TPW.3}
\]
No old label or weight is removed from this established union.

\paragraph{The additional growing family.}
Let $\Xi_7$ be the six Dirichlet characters modulo seven, using their
associated primitive characters; the trivial one has conductor one.
For $\xi\in\Xi_7$ set
\[
 M(\xi)=
 \begin{cases}
 14,&\xi=1,\\
 17,&\operatorname{ord}(\xi)=2,\\
 19,&\operatorname{ord}(\xi)=3\text{ or }6.
 \end{cases}
 \tag{TPW.4}
\]
For an integer $h\geq2$ define
\[
 \begin{aligned}
 \mathscr A_h=\{\,&(\delta\xi)|_{G_K}:\ \xi\in\Xi_7,\
       \delta\text{ primitive of exact conductor }2^c,\\
       &M(\xi)\leq c\leq h+2\,\},\\
 Z^{\rm high2}(u)&=\{(p,\eta,e):\eta\in\mathscr A_{h(u)},\ 1\leq e<T\},\\
 Z^{\rm unif2}(u)&=Z^{\rm qref}(u)\sqcup Z^{\rm high2}(u).
 \end{aligned}
 \tag{TPW.5}
\]
Both parities of $\delta$ occur.  These labels have trivial cubic and
trivial auxiliary $q,d$ factors.  Their finite factor at $269$ is
trivial.  In particular the Teichmuller compensation and near-one lift
allowed in AC-G173 are not new finite host labels in (TPW.5).

For positive integer $T$, put
\[
 \begin{aligned}
 N_T&=T-1,&B_T&=\lfloor(T-1)/268\rfloor,& A_T&=\lfloor T/67\rfloor,\\
 R_T&=\sum_{r\in\mathcal E}\left\lfloor\frac{T+267-r}{268}\right\rfloor,\\
 S_T&=\sum_{\sigma=\pm1}\sum_{b=0}^1\sum_{r\in\mathcal F_{\sigma,b}}
                         \left\lfloor\frac{T+267-r}{268}\right\rfloor,\\
 B(h;M)&=\max\{0,2^{h+1}-2^{M-2}\},\\
 C(h)&=B(h;14)+B(h;17)+4B(h;19).
 \end{aligned}
 \tag{TPW.6}
\]
$B_T$ and $B(h;M)$ are different quantities.
Let $I_{399}^{\rm unif2}$ omit precisely $Z^{\rm unif2}$ from (TC.1).
Its immediate predecessor is $I_{399}^{\rm estquartreflection}$.

\begin{theorem}[AC-G174-window: one growing-family reduction on the original window]
\label{AC-G174-window}
Every triple in (TPW.3) and (TPW.5) is an actual character, coefficient
prime and weight whose original strict group is zero.
The new family is disjoint from every previously omitted term, and
\[
 \begin{aligned}
 |\mathscr A_h|&=C(h),\\
 |Z^{\rm unif2}\setminus Z^{\rm qref}|&=C(h)N_T,\\
 |Z^{\rm unif2}_{269}|&=(8+C(h))N_T-B_T+12A_T-R_T-S_T,\\
 |Z^{\rm unif2}_3|&=N_T.
 \end{aligned}
 \tag{TPW.7}
\]
For $h\geq17$, $C(h)=12\cdot2^h-561152$.
For $h\geq12$ and $T\geq2$ the inclusion of omitted sets is proper;
these inequalities hold eventually on every actual page (TPW.1).
This is a proper enlargement of the proved set, not an assertion that
the numerical nonnegative sum must strictly decrease.

Let $E_u$ denote the normalized positive-index cost of the entire union.
Then in nonnegative extended sums,
\[
 \begin{aligned}
 I_{399}&=I_{399}^{\rm unif2}+E_u,\\
 I_{399}^{\rm unif2}&\leq I_{399}^{\rm estquartreflection}\leq I_{399},\\
 0\leq E_u&\leq8D_u^2T^2\log(8TD_u)
                 +\tfrac12D_u\mathcal R_uT^2+2D_ur_uc_0T\\
 &=O_\varepsilon\bigl(u^{802+16\varepsilon}(\log u)^2\bigr).
 \end{aligned}
 \tag{TPW.8}
\]
There is one full (FL.4) cost for the union.  It is not multiplied by
$C(h)$, and no individual derived index is asserted zero.
The new positive-relative-slack premise implies the original full
common-window premise on the same $W,J$, datum, page, pair and original
coefficients, after only increasing the threshold and halving slack.
Every older premise implies the new premise by nonnegativity.
The actual $R.t$, all simultaneous field quantifiers and the retained
mixed realization $\mathcal R_{\rm cub}$ are unchanged.
\end{theorem}

\begin{proof}
\proofstep{1. Actual characters, original coefficients and strict kernels.}
The host contains $\mathbb Q(\zeta_{2^{h+2}},\zeta_7)$.
Every label in $\mathscr A_h$ therefore factors through $G_u$.
The intersection of this rational cyclotomic field with $K$ is
$\mathbb Q$: the former is unramified at three and $K/\mathbb Q$
ramifies there.  Distinct pairs $(\delta,\xi)$ stay distinct after
restriction to $G_K$, since the two and seven conductor factors are
independent and their ratio cannot be $\chi_{-3}$.
A primitive character of conductor $2^c$, $c\geq3$, has order dividing
$2^{c-2}$.  Thus each product has order dividing $3\cdot2^h$, exactly
the exponent of (FC.1).  Its values lie in the same original $O_p$.

For each new label $\eta=(\delta\xi)|_{G_K}$ the actual coefficient is
\[
 \begin{aligned}
 A_{\eta,e}
  &=(O_p\otimes_{\mathbb Z_p}\mathbb Q_p/\mathbb Z_p)(\eta)(-e)\\
  &=(E_p/O_p)(\delta\xi c_p^{-e}).
 \end{aligned}
 \tag{TPW.9}
\]
The equality is the canonical one for the finite free $\mathbb Z_p$
module $O_p$; neither its rank nor its lattice is changed.
Apply (G173.17) with this $O_p$, $s=0$ and $\kappa=1$.
All inequalities on $c$ are exactly (TPW.4)--(TPW.5).
It proves strict zero for every weight, including residue zero,
on the arbitrary original coefficient ring.
In particular this uses G173's full $K_p$ Frobenius argument in
residue zero, not merely a residual class-group calculation.

Inflation from $G_{K,\Sigma_p}$ to $G_K$ is injective because the
coefficient is unramified outside $\Sigma_p$.
An inflated (FC.2) class kills inertia outside that set, satisfies
the same away-inertia zero conditions inside it, and is zero in
full local $H^1$ at $p$.  It is therefore a class of (G173.16) and
vanishes.  Equivalently, when comparing smaller ramification sets,
cohomologically zero inertia restrictions for an unramified coefficient
are zero cocycles; the cocycle identity kills their conjugates and
closed normal closure.  All comparisons preserve the original module
and its full local zero condition.
Every old strict zero in (TPW.3) is retained from AC-G171-window.

\proofstep{2. Disjointness and the exact growing count.}
All old pure finite labels have two-primary conductor at most sixteen.
Every new label has two-primary conductor at least $2^{14}$.
Their ratio cannot become the conductor-three character on passing to
$G_K$, so no new label equals an old pure one.
The nontrivial-cubic old labels are distinct as well: $F_u/\mathbb Q$
is abelian, whereas $K(t)/\mathbb Q$ is the nonabelian $S_3$ splitting
field; hence $F_u\cap K(t)=K$ and the cubic factor is independent.
The coefficient-three row is separate.  This proves the disjoint union.

For $c\geq3$, the number of primitive Dirichlet characters of conductor
$2^c$ is $2^{c-2}$: subtract the $2^{c-2}$ characters factoring through
modulus $2^{c-1}$ from all $2^{c-1}$ characters modulo $2^c$.
Consequently
\[
 \sum_{c=M}^{h+2}2^{c-2}=2^{h+1}-2^{M-2}
 \quad(M\leq h+2),
 \tag{TPW.10}
\]
with zero if the range is empty.  There is one trivial character modulo
seven, one quadratic one, two of order three and two of order six.
This proves $C(h)$, its large-$h$ formula, and the new count $C(h)N_T$.
The old coefficient-$269$ count is $8N_T-B_T+12A_T-R_T-S_T$,
and the old coefficient-three count is $N_T$.
Adding the disjoint new set gives (TPW.7).
At $h\geq12$ at least the conductor-$2^{14}$ characters with $\xi=1$
occur; together with $e=1$ this proves proper inclusion for $T\geq2$.

\proofstep{3. The single full-union local cost.}
Keep the original codifferent dual and normalizing rank $b_\ell$.
For every omitted triple, strict zero and (EU.2) give
\[
 \begin{aligned}
 \frac{\log\#k_{O_\ell}}{b_\ell}\max(0,i_{\ell,\eta,e})
 &\leq\frac{\log\#H^2(C_{\ell,\eta,e})}{b_\ell}\\
 &=\frac{\left|\log\#H^2(C_{\ell,\eta,e})
                  -\log\#\mathcal S_{\ell,\eta}(e)\right|}{b_\ell}.
 \end{aligned}
 \tag{TPW.11}
\]
For each weight the entire enlarged set is a subsum of the nonnegative
all-prime, all-character (FL.4) error.  That bound already includes
every character of $G_u$; increasing the size of our subset does not
multiply it by the number of new labels.
Sum its bound $16eD_u^2\log(8eD_u)+D_ue\mathcal R_u+2D_ur_uc_0$
once over $1\leq e<T$, using $\sum e\leq T^2/2$, to obtain (TPW.8).
The extended-sum equality follows by nonnegative additivity, without
subtracting infinities.
The unchanged $D_u\leq288u^{2+8\varepsilon}$, $r_u\leq12$,
$\mathcal R_u=O_\varepsilon(u^{\varepsilon/5}(\log u)^2)$ and
$T\leq u^{399}$ give the exponent $802+16\varepsilon$.
It is little-$o$ of $\exp(u^\varepsilon/2)$ for every fixed positive
$\varepsilon$, uniformly in the prescribed pair.

\proofstep{4. The same prescribed window and actual rank.}
Fix the retained finite datum, page and $\varepsilon$.
Suppose the new premise supplies $J\geq1$, $X_0$ and $s\in(0,1)$,
independent of $X,q,d,u,\ell,\eta,e$, such that every $X\geq X_0$
and every prescribed $\operatorname{WinPair}(\varepsilon,X^{1/J},q,d)$
have one $W$ with $X^{1/J}\leq W$, $W^2\leq X$, and
\[
 I_{399}^{\rm unif2}(H_u,u)\leq(1-s)\exp(u^\varepsilon/2)
                       \qquad(W\leq u\leq W^2).
 \tag{TPW.12}
\]
Choose $U_0\geq2$ so the single error in (TPW.8) is at most
$(s/2)\exp(u^\varepsilon/2)$ for $u\geq U_0$.
Replacing $X_0$ by $\max(X_0,U_0^J)$ leaves exactly the supplied $W$
and gives the original full premise with slack $s/2$ there.
The same proof works for uniform little-$o$ premises and for every
fixed finite simultaneous family, by taking a maximum of its uniform
error constants on that same window.
Conversely, nonnegativity gives the new premise from every older one
with unchanged parameters.  No new selection or intersection of windows
is used, and no stronger compatibility condition is imposed.
The source rank $R.t$, all field/weight/character/coefficient quantifiers,
and the unchanged mixed theta/Selmer and analytic realization inputs
are exactly those of AC-K35.
\end{proof}

\paragraph{Remaining scope and finite-cover dictionary.}
Only the stated coefficient-$269$ labels are added.
All unlisted coefficient primes, the low-conductor exceptions retained
in (TPW.3), new $q,d$ twists and all remaining nontrivial-cubic labels
stay in the unknown sum.  The theorem does not classify every vanishing
character or prove that an uncovered term is nonzero.
It gives no full numerical bound for the remainder,
no mixed realization $\mathcal R_{\rm cub}$ and no H04.4 conclusion
without the unchanged retained inputs.
The new characters are rank-one Artin factors on actual finite
cyclotomic covers; their Tate tensor and original integral cohomology
are the existing AC-K25/26 dictionary.  The audited SGA~1,
Expos\'e~V, \S\S4,7 and Stacks~03PK scopes are retained.
No new motive, Hodge comparison, archive scan or Lean execution is claimed.
